% !TeX program = lualatex
% Compile twice: lualatex 154.tex
% All references and prose are embedded; no BibTeX database or external inputs.
\documentclass[11pt,a4paper]{article}
\usepackage[margin=24mm,headheight=14pt]{geometry}
\usepackage{fontspec}
\setmainfont{Latin Modern Roman}
\setsansfont{Latin Modern Sans}
\setmonofont{Latin Modern Mono}
\usepackage{amsmath,amssymb,mathtools}
\usepackage{unicode-math}
\setmathfont{Latin Modern Math}
\usepackage{microtype}
\usepackage{enumitem,xurl,needspace}
\usepackage[dvipsnames]{xcolor}
\usepackage{hyperref}
\hypersetup{unicode=true,colorlinks=true,linkcolor=MidnightBlue,urlcolor=MidnightBlue,
 pdftitle={500 Mathematical Problem Statements},pdfauthor={Source-annotated compilation},bookmarksnumbered=true}
\usepackage{bookmark}
\usepackage{tocloft}
\setlength{\cftsecnumwidth}{3em}
\cftsetpnumwidth{2.5em}
\cftsetrmarg{4em}
\usepackage{titlesec}
\titleformat{\section}{\Large\bfseries\raggedright}{\thesection}{0.7em}{}
\usepackage{fancyhdr}
\pagestyle{fancy}\fancyhf{}
\fancyhead[L]{\small\sffamily Mathematical problem statements}
\fancyhead[R]{\small Source edition 22}
\fancyfoot[C]{\thepage}
\setlength{\parindent}{0pt}
\setlength{\parskip}{5pt plus 1pt}
\setlength{\emergencystretch}{3em}
\setcounter{tocdepth}{1}
\setcounter{secnumdepth}{1}
\setlist[itemize]{leftmargin=3em,itemsep=5pt,topsep=4pt,parsep=0pt}
\allowdisplaybreaks
\newcommand{\N}{\mathbb{N}}
\newcommand{\Z}{\mathbb{Z}}
\newcommand{\Q}{\mathbb{Q}}
\newcommand{\R}{\mathbb{R}}
\newcommand{\C}{\mathbb{C}}
\newcommand{\F}{\mathbb{F}}
\newcommand{\A}{\mathbb{A}}
\newcommand{\PP}{\mathbb P}
\newcommand{\E}{\mathbb{E}}
\newcommand{\eps}{\varepsilon}
\newcommand{\dd}{\,\mathrm d}
\newcommand{\sref}[2]{\hyperlink{source:#1:#2}{[S#2]}}
\newcommand{\eref}[2]{\hyperlink{extra:#1:#2}{[E#2]}}
% Turn the next switch off for a shorter reading copy. The quotations preserve
% every exactTarget field, including qualifications omitted from a short summary.
\newif\ifshowcatalogue
\showcataloguetrue
\newcommand{\cataloguescope}[1]{\ifshowcatalogue
 \paragraph{Exact scope in the supplied catalogue.}
 \begin{quote}\small #1\end{quote}\fi}

% Standalone research-notebook wrapper; the source preamble above is retained.
\numberwithin{equation}{section}
\hypersetup{pdftitle={Research attempt -- problem 154},
 pdfauthor={Research notebook; original compilation retained}}
\fancyhead[L]{\small\sffamily Research notebook}
\fancyhead[R]{\small\sffamily Problem 154 | 27 September 2026}
\begin{document}
\setcounter{section}{153}
% ================================================================
% problemId: problem.top500-omission-w045049-hyperlinearity-of-every-group
% source releaseRank: 154; source releaseStatus: open
% exactTarget SHA-256: 39a985acb5b38b04747a5c888a3ee595a0636175e9f8a616841c8c02b9447170
\clearpage
\section{\texorpdfstring{Hyperlinearity of Every Group}{Hyperlinearity of Every Group}}\label{154}
\subsection{Definitions and mathematical statement}
Let $U(d)=\{U\in M_d(\C):U^*U=I\}$ and
$\|A\|_{2,d}=(d^{-1}\operatorname{Tr}(A^*A))^{1/2}$. A group $G$ is hyperlinear in the selected convention if, for every finite $F\subseteq G$ and $\eps>0$, there is a map $\phi:F\to U(d)$ for some $d$ such that
\[
 \|\phi(g)\phi(h)-\phi(gh)\|_{2,d}<\eps,
 \qquad \|\phi(g)-I\|_{2,d}>\sqrt2-\eps\quad(g\ne e).
\]
The first condition applies only when $g,h,gh\in F$, and the second only to elements of $F$. Dimension may depend on $F$ and $\eps$; exact homomorphisms are not required. Determine whether every group admits these finite approximate matrix models.
\par\smallskip\textit{Formulation and concept references:} \sref{154}{1}.
\cataloguescope{Determine whether every group G is hyperlinear. Precisely, for each finite F contained in G and each epsilon>0, must there be d>=1 and phi:F->U(d) such that ||phi(g)phi(h)-phi(gh)||\_(2,d)<epsilon whenever g,h,gh belong to F, and ||phi(g)-I||\_(2,d)>sqrt(2)-epsilon for every nonidentity g in F? Here ||X||\_(2,d)=(Tr(X*X)/d)\textasciicircum{}(1/2) is the normalized Hilbert-Schmidt norm. A negative resolution is a group failing this condition.}
\subsection{Short English statement}
Can every finite portion of every group be modeled by unitary matrices that nearly preserve multiplication while keeping nonidentity elements detectably away from the identity?
% BEGIN RESEARCH ADDITION ATTEMPT-154
\subsection{Research attempt}

\paragraph{Current frontier.}
Dogon's Theorems 1.3 and 1.6 give non-hyperlinearity of suitable nonsplit
central extensions under a flexible Hilbert--Schmidt stability, or weaker
ucp-stability, hypothesis. Property (T) alone is not that hypothesis
\sref{154}{1}. A manuscript by Jihao Liu, in its consulted 21 September 2026 version, claims an
explicit non-hyperlinear wreath product, using internality and normalization
of commutants in matrix ultraproducts \eref{154}{1}. \textbf{Editorial status
note:} the original catalogue text is retained. That manuscript is a current
proof claim, not an independently verified theorem in this notebook. The
calculations below check a specific obstruction and its compatibility with
the \emph{printed} approximation convention; they do not audit its full
analytic argument.

\paragraph{Attack route.}
Three routes have rather different missing inputs. The central-extension
route needs a stability theorem strong enough to invoke Dogon. A route
through nonembeddable operator algebras needs the obstruction to concern a
\emph{group's canonical trace}, rather than just some trace on an algebra.
The route pursued here instead tries to force two distinct lamps in a
permutation wreath product to have the same image in every matrix
ultraproduct. Its precise missing input is normalization of a subgroup
commutant. This gives a useful finite-error question, even without assuming
that the recent claim has been verified.

\paragraph{Attempt: an explicit group and a nontrivial word.}
Use the construction in \eref{154}{1}, Theorem 1.1:
\[
 R_+=\F_2[x_1,x_2,x_3],\qquad
 R=\F_2[x_1^{\pm1},x_2^{\pm1},x_3^{\pm1}],\qquad
 \Gamma=\operatorname{EL}_3(R_+),
\]
\[
 G=\operatorname{EL}_3(R)\rtimes_\alpha\operatorname{SL}_3(\Z),\qquad
 W=\left(\bigoplus_{G/\Gamma}C_2\right)\rtimes G.
\]
Here $\operatorname{EL}_3(B)$ is generated by $I_3+bE_{ij}$, $i\ne j$,
and $\alpha_A(x^v)=x^{Av}$ for column vectors $v\in\Z^3$.
Write $a_q$ for the nonidentity lamp at the left coset $q$. Set
\[
 t=(I_3,I_3+E_{12}),\qquad
 \gamma=(I_3+x_2E_{12},I_3),\qquad
 x=a_{t\Gamma},\quad y=\gamma x\gamma^{-1},\quad b=yx.
\]
The following verifications are supplied here rather than inferred from the
paper's claimed conclusion. The inverse monomial shear sends $x_2$ to
$x_1^{-1}x_2$, so
\[
 t^{-1}\gamma t=(I_3+x_1^{-1}x_2E_{12},I_3)\notin\Gamma.
\]
Every matrix in $\operatorname{EL}_3(R_+)$ has polynomial entries, which
justifies the last exclusion. Thus $\gamma t\Gamma\ne t\Gamma$. The lamps
at these two cosets are independent involutions, so $x^2=y^2=e$ and
$b\ne e$. Notice also that $t\Gamma t^{-1}\subseteq\Gamma$, since the
forward shear sends polynomial monomials to polynomial monomials.

Let $M$ be a tracial matrix ultraproduct and $\rho:W\to U(M)$ a
homomorphism. Put $D=\rho(\Gamma)'\cap M$.
We isolate the following input, without claiming to establish it:
\begin{equation}\label{eq154-normalize}
 \rho(t)D\rho(t)^*=D
 \quad\text{for every such }M\text{ and }\rho.
\end{equation}
\textbf{Conditional lamp lemma.} Equation~\eqref{eq154-normalize} implies
$\rho(b)=I$ for every $\rho$.

\textit{Proof.} The lamp $a_\Gamma$ commutes with $\Gamma$, since this
subgroup fixes its coset. Hence $\rho(a_\Gamma)\in D$. Normalization gives
$X:=\rho(x)=\rho(t)\rho(a_\Gamma)\rho(t)^*\in D$. In particular $X$
commutes with $\rho(\gamma)$. Consequently
$\rho(y)=X$, and $\rho(b)=X^2=I$. This uses an actual equality of
commutants, not merely a one-sided inclusion. \hfill$\square$

The recent manuscript asserts centralizer results intended to supply this
normalization in its wreath-product application; see its Theorems 6.7, 7.4,
and 8.3 \eref{154}{1}. We keep \eqref{eq154-normalize} as an unproved input
in the present argument.

\paragraph{Attempt: a dimension-free finite-error obstruction.}
The lamp argument admits an explicit estimate that does not require
ultraproduct terminology. Let $F$ contain
\[
 e,\ \gamma,\ \gamma^{-1},\ x,\ \gamma x,\ y,\ b,
\]
and suppose $\phi:F\to U(d)$ satisfies the printed multiplicative
inequalities with tolerance $\eps$. Put $H=\phi(\gamma)$ and
$X=\phi(x)$, and use normalized Hilbert--Schmidt norms throughout.
All the products needed below have their factors and product in $F$.
From $e^2=e$ and $\gamma\gamma^{-1}=e$ we obtain
\[
 \|\phi(e)-I\|_2<\eps,\qquad
 \|\phi(\gamma^{-1})-H^*\|_2<2\eps.
\]
Using $(\gamma x)\gamma^{-1}=y$, followed by $x^2=e$ and $yx=b$, gives
\[
 \|\phi(y)-HXH^*\|_2<4\eps,\qquad
 \|X^2-I\|_2<2\eps,
\]
and therefore
\begin{equation}\label{eq154-error}
 \|\phi(b)-I\|_2
 <7\eps+\|HX-XH\|_2
 \le 7\eps+2\operatorname{dist}_2(X,\{H\}').
\end{equation}
For the last inequality, subtract any matrix commuting with $H$ and use
unitary invariance of the norm. No positivity or unitarity is required of
that comparison matrix. The separation requirement for $b\ne e$ now
forces
\begin{equation}\label{eq154-distance}
 \operatorname{dist}_2(X,\{H\}')>
 \frac{\sqrt2-8\eps}{2}.
\end{equation}
In particular, any proposed uniform repair argument must overcome an
order-one distance, not just improve a word-error constant. Equations
\eqref{eq154-error}--\eqref{eq154-distance} are unconditional finite-matrix
lemmas proved here; they do not assert that the distance is small.

\paragraph{Attempt: the repair estimate that really is available.}
Suppose $\pi:\Lambda\to U(d)$ is an \emph{exact} representation and a finite
set $S\subset\Lambda$ has an adjoint spectral gap $\kappa>0$, in the
following precise sense:
\[
 \sum_{s\in S}\|\pi(s)V\pi(s)^*-V\|_2^2
 \ge\kappa^2\|V\|_2^2
 \quad\text{whenever }V\perp\pi(\Lambda)'.
\]
Let $E$ be orthogonal projection onto the commutant in the Hilbert space
$M_d(\C)$ with normalized trace inner product. Applying the gap to $A-E(A)$
proves
\begin{equation}\label{eq154-gap}
 \|A-E(A)\|_2\le\kappa^{-1}
 \left(\sum_{s\in S}\|[A,\pi(s)]\|_2^2\right)^{1/2}.
\end{equation}
Indeed, $E(A)$ is fixed by conjugation and the orthogonal complement is
invariant. For unitary $A$, this projection is also the trace-preserving
conditional expectation onto the finite-dimensional commutant; hence
$Y=E(A)$ is a contraction. Extend the polar partial isometry of $Y$ to a
unitary $U$ in each matrix block of the commutant. Such an extension exists
because initial and final ranks agree in every square block. Then
\[
 \|Y-U\|_2^2=\operatorname{tr}_d(1-|Y|)^2
 \le1-\|Y\|_2^2=\|A-Y\|_2^2,
\]
so $\|A-U\|_2\le2\|A-E(A)\|_2$.
This proves a concrete commuting-unitary repair theorem under the stated
\emph{coordinate-level} gap hypothesis.

It does not prove \eqref{eq154-normalize}. Lifts of an ultraproduct
representation need only be approximate representations in each coordinate.
The spectral-gap inequality above has not been proved for those lifts, and
replacing them by exact representations would import a stability premise of
the kind isolated in \sref{154}{1}.

\paragraph{Stress test.}
First, almost commuting need not mean close to the exact coordinate
commutant. Take
\[
 U_n=\begin{pmatrix}1&0\\0&e^{i/n}\end{pmatrix},\qquad
 V=\begin{pmatrix}0&1\\1&0\end{pmatrix}.
\]
Then $\|[U_n,V]\|_{2,2}=|1-e^{i/n}|\to0$, whereas the commutant of $U_n$
is diagonal and $\operatorname{dist}_{2,2}(V,\{U_n\}')=1$. The limiting
unitary is $I$, whose commutant is much larger. This is not a counterexample
to a property-(T) theorem; it disproves an unjustified passage from limiting
commutation to exact coordinate commutants.

Second, a trace-preserving one-sided inclusion need not be equality in an
infinite finite von Neumann algebra. For example, inside the bilateral
infinite tensor product $\overline\bigotimes_{k\in\Z}M_2$, the shift sends
$D_+=\overline\bigotimes_{k\ge0}M_2$ properly into itself. In the crossed
product by the shift this map is implemented by a unitary and preserves the
trace. Thus a finite-dimensional dimension-counting argument cannot simply
be transferred to the ultraproduct. This example is not an asserted model
of the specified property-(T) construction.

Finally, we check the exact convention in the catalogue. Assume that the
printed finite approximations exist for the countable group $W$. Choose
increasing finite sets exhausting $W$, and approximations with
$\eps_n\to0$. Extend each map arbitrarily by $I$ outside its domain. The
classes of these bounded unitary sequences in any nonprincipal tracial
matrix ultraproduct give an exact homomorphism $\rho$. The elementary
estimate on $\phi(e)$ above gives $\rho(e)=I$; no additional normalization
was assumed. Moreover
\[
 \|\rho(b)-I\|_2=\lim_{n\to\omega}\|\phi_n(b)-I\|_2\ge\sqrt2.
\]
This contradicts the conditional lamp lemma. Thus
\eqref{eq154-normalize}, if justified, really would disprove the exact
statement here, not merely a stronger pairwise-separation variant.

\paragraph{Outcome.}
\textbf{Outcome: Conditional reduction.}

The supplied argument proves nontriviality of the candidate word, the
finite-error obstruction \eqref{eq154-error}, and the implication from
commutant normalization to failure of the catalogue's approximations. It
also proves a repair lemma for exact representations and exhibits why its
naive use on approximate representations fails. These deductions are not a
new construction of the underlying group, and no priority is claimed for
the conditional lamp mechanism.

The unresolved step in this notebook is the uniform normalization input
\eqref{eq154-normalize}; the recent manuscript claims to provide it, but
its full proof has not been independently verified here. Establishing that
input would produce a negative answer and would also rule out soficity of
this group, since permutation-matrix approximations are a special source
of hyperlinear approximations \sref{154}{1}, Introduction. The weaker
finite-matrix estimates alone do not supply a counterexample.
% END RESEARCH ADDITION ATTEMPT-154
\subsection{Sources}
\begin{itemize}\raggedright
\item[\textnormal{[S1]}]\hypertarget{source:154:1}{} A. Dogon, Flexible Hilbert-Schmidt Stability versus Hyperlinearity for Property (T) Groups, arXiv:2211.10492v3 (25 August 2023).
\url{https://arxiv.org/html/2211.10492v3}.
{\small\textit{Catalogue formulation source.}}
% BEGIN RESEARCH ADDITION SOURCES-154
\item[\textnormal{[E1]}]\hypertarget{extra:154:1}{} Jihao Liu,
\emph{Nonhyperlinear groups exist}, author-hosted manuscript, first posted
20 September 2026; consulted PDF dated 21 September 2026, Theorems 1.1--1.3, 6.7, 7.4, and 8.3.
\url{https://jihaoliu.org/ai-results/nonhyperlinear-groups-exist-2026-09-20.pdf}.
{\small\textit{Consulted 27 September 2026. A proof claim; not treated as
independently verified. The explicit group is attributed to this source.}}
% END RESEARCH ADDITION SOURCES-154
\end{itemize}

\end{document}
